\chapter{LLD Case Studies}
\chapterepigraph{An LLD interview is a conversation, not a coding sprint.
You are judged on how you clarify vague requirements, carve the problem
into classes with clean responsibilities, choose patterns with
justification, and leave room for the features the interviewer will ask
for next. Each case study here follows that arc: clarify, identify
entities, design, and name the patterns that keep the design flexible.}

\section{How to Approach Any LLD Problem}
\label{sec:cases-intro}

Before writing a single class, run the same four steps every time. They
turn a scary open-ended prompt into a structured design.

\begin{ideabox}[The Four-Step LLD Method]
\begin{enumerate}
  \item \textbf{Clarify \& scope.} Ask about requirements and
        constraints; state your assumptions. Nail down what is in and out
        of scope \emph{before} designing.
  \item \textbf{Identify entities.} Extract the nouns -- they become
        classes; the verbs become methods. Assign each class one clear
        responsibility (SRP).
  \item \textbf{Design relationships \& patterns.} Decide how classes
        relate (composition over inheritance) and where a known pattern
        fits (Strategy for varying rules, Factory for creation, Observer
        for notifications, State for lifecycles).
  \item \textbf{Walk a scenario \& extend.} Trace one end-to-end flow,
        then show how a likely new requirement slots in without rewrites.
\end{enumerate}
\end{ideabox}

\begin{patternbox}[Patterns That Recur Across Case Studies]
\begin{itemize}
  \item \textbf{Pricing/matching/eligibility rules that vary} --
        \textbf{Strategy}.
  \item \textbf{Creating objects by type} -- \textbf{Factory}.
  \item \textbf{Lifecycle with modes} (order, game, machine) --
        \textbf{State}.
  \item \textbf{``Notify these when that happens''} -- \textbf{Observer}.
  \item \textbf{One shared resource} (config, id generator, in-memory
        store) -- \textbf{Singleton} (used judiciously).
  \item \textbf{Complex object assembly} -- \textbf{Builder}.
\end{itemize}
\end{patternbox}

The case studies progress from the classic ``design a parking lot'' to
notification, logging, and rate-limiting systems. Each names the patterns
from Chapter~2 it relies on, so the two chapters reinforce each other.

\newpage

\section{Design a Parking Lot}
\label{sec:case-parking-lot}

\problemstatement{Design a parking lot with multiple levels, multiple spot
sizes (motorcycle, compact, large), vehicle entry/exit, ticketing, and fee
calculation. It should assign the smallest spot a vehicle fits in and
compute the fee on exit.}

\begin{patternbox}[Clarify Before Designing]
Ask: How many levels and spot types? Can a large vehicle occupy multiple
spots? Is pricing flat, hourly, or tiered? Multiple entry/exit gates?
Payment methods? For this design assume: several levels, three spot sizes,
one spot per vehicle, hourly pricing that may vary by vehicle type, and
multiple gates.
\end{patternbox}

\subsection*{Identify the entities}

The nouns give the classes: \texttt{ParkingLot} (owns levels),
\texttt{Level} (owns spots), \texttt{ParkingSpot} (has a size and
occupancy), \texttt{Vehicle} (abstract; \texttt{Motorcycle},
\texttt{Car}, \texttt{Truck}), \texttt{Ticket} (issued at entry),
and a \texttt{FeeStrategy} for pricing. The verbs -- park, unpark, pay --
become methods.

\begin{ideabox}[Patterns That Keep It Flexible]
\textbf{Strategy} for \texttt{FeeStrategy} so pricing rules change without
touching parking logic. \textbf{Factory} to create the right
\texttt{Vehicle}/spot type. A \textbf{Singleton} (or injected single
instance) \texttt{ParkingLot} as the entry point. Spot assignment is a
policy method, easy to swap (nearest, smallest-fit).
\end{ideabox}

\subsection*{Class design (C++ skeleton)}

\begin{lstlisting}
#include <bits/stdc++.h>
using namespace std;

enum class SpotSize { Motorcycle, Compact, Large };

class Vehicle {
public:
    string plate;
    SpotSize required;
    Vehicle(string p, SpotSize s) : plate(move(p)), required(s) {}
    virtual ~Vehicle() = default;
};
struct Car : Vehicle { Car(string p) : Vehicle(move(p), SpotSize::Compact) {} };
struct Bike : Vehicle { Bike(string p) : Vehicle(move(p), SpotSize::Motorcycle) {} };

struct FeeStrategy {                            // Strategy: pricing varies
    virtual double fee(int hours) const = 0;
    virtual ~FeeStrategy() = default;
};
struct FlatHourly : FeeStrategy {
    double rate; FlatHourly(double r) : rate(r) {}
    double fee(int hours) const override { return rate * hours; }
};

class ParkingSpot {
public:
    int id; SpotSize size; Vehicle* occupant = nullptr;
    ParkingSpot(int i, SpotSize s) : id(i), size(s) {}
    bool canFit(const Vehicle& v) const {
        return occupant == nullptr && size >= v.required;   // smallest-fit rule
    }
};

struct Ticket { int spotId; string plate; time_t entry; };

class Level {
    vector<ParkingSpot> spots;
public:
    Level(int n) { for (int i = 0; i < n; i++) spots.emplace_back(i, SpotSize::Compact); }
    ParkingSpot* findSpot(const Vehicle& v) {
        for (auto& s : spots) if (s.canFit(v)) return &s;   // assignment policy
        return nullptr;
    }
};

class ParkingLot {
    vector<Level> levels;
    unique_ptr<FeeStrategy> pricing;
public:
    ParkingLot(unique_ptr<FeeStrategy> p) : pricing(move(p)) {}
    void addLevel(int spots) { levels.emplace_back(spots); }

    Ticket* park(Vehicle& v) {
        for (auto& level : levels) {
            if (auto* spot = level.findSpot(v)) {
                spot->occupant = &v;
                return new Ticket{ spot->id, v.plate, time(nullptr) };
            }
        }
        return nullptr;                         // lot full
    }
    double unpark(Ticket* t, int hours) { return pricing->fee(hours); }
};

int main() {
    ParkingLot lot(make_unique<FlatHourly>(20.0));
    lot.addLevel(2);
    Car car("KA-01-1234");
    Ticket* t = lot.park(car);
    cout << (t ? "Parked, fee for 3h = " : "Full ") ;
    if (t) cout << lot.unpark(t, 3) << "\n";
    return 0;
}
\end{lstlisting}

\begin{takeawaybox}
Parking lot is the canonical LLD warm-up. Score points by: clarifying
scope first, giving each class one job (SRP), isolating pricing behind a
\textbf{Strategy}, and making spot-assignment a swappable policy. Likely
follow-ups: multiple gates (concurrency, locking a spot atomically),
payment methods (another Strategy), and reservations.
\end{takeawaybox}

\section{Design an Elevator System}
\label{sec:case-elevator}

\problemstatement{Design the control system for a building with multiple
elevators. Handle external hall calls (up/down on a floor) and internal
car requests (go to floor N), decide which elevator serves a request, and
move each car efficiently.}

\begin{patternbox}[Clarify Before Designing]
Ask: How many elevators and floors? What scheduling goal -- minimise wait,
minimise travel, or fairness? Are there express/service elevators? For
this design assume: $N$ elevators, a central dispatcher, a
nearest-suitable-car scheduling policy that is swappable, and each car
running a SCAN (elevator algorithm) sweep.
\end{patternbox}

\subsection*{Identify the entities}

\texttt{Elevator} (a car with current floor, direction, and its set of
target floors), \texttt{Request} (hall or car call), \texttt{Direction}
(up/down/idle), a \texttt{Dispatcher} that assigns requests to cars via a
\texttt{SchedulingStrategy}, and a \texttt{Building} owning the elevators.
Each elevator behaves as a small state machine: idle, moving up, moving
down.

\begin{ideabox}[Patterns That Keep It Flexible]
\textbf{Strategy} for \texttt{SchedulingStrategy} (nearest car, least
load, zoning) so dispatch policy is a plug-in. \textbf{State} for each
elevator's motion (idle/up/down) to keep the movement logic clean.
Optionally \textbf{Observer} so displays update as a car moves. The car's
own sweep is the classic SCAN scheduling of its pending stops.
\end{ideabox}

\subsection*{Class design (C++ skeleton)}

\begin{lstlisting}
#include <bits/stdc++.h>
using namespace std;

enum class Direction { Up, Down, Idle };

class Elevator {
public:
    int id, floor = 0;
    Direction dir = Direction::Idle;
    set<int> targets;                           // pending stops

    void addTarget(int f) { targets.insert(f); }
    int distanceTo(int f) const { return abs(floor - f); }

    void step() {                               // one SCAN move
        if (targets.empty()) { dir = Direction::Idle; return; }
        int next = (dir == Direction::Down)
                     ? *prev(targets.upper_bound(floor))
                     : *targets.lower_bound(floor);
        if (next > floor) { floor++; dir = Direction::Up; }
        else if (next < floor) { floor--; dir = Direction::Down; }
        else { targets.erase(floor); cout << "Car " << id << " stop @ " << floor << "\n"; }
    }
};

struct SchedulingStrategy {                     // Strategy: which car serves
    virtual Elevator* choose(vector<Elevator>& cars, int floor) = 0;
    virtual ~SchedulingStrategy() = default;
};
struct NearestCar : SchedulingStrategy {
    Elevator* choose(vector<Elevator>& cars, int floor) override {
        Elevator* best = nullptr;
        for (auto& c : cars)
            if (!best || c.distanceTo(floor) < best->distanceTo(floor)) best = &c;
        return best;
    }
};

class Dispatcher {
    vector<Elevator> cars;
    unique_ptr<SchedulingStrategy> policy;
public:
    Dispatcher(int n, unique_ptr<SchedulingStrategy> p) : policy(move(p)) {
        for (int i = 0; i < n; i++) cars.push_back({i});
    }
    void hallCall(int floor) {
        Elevator* car = policy->choose(cars, floor);
        car->addTarget(floor);
        cout << "Assigned car " << car->id << " to floor " << floor << "\n";
    }
    void tick() { for (auto& c : cars) c.step(); }
};

int main() {
    Dispatcher d(2, make_unique<NearestCar>());
    d.hallCall(5);
    for (int i = 0; i < 6; i++) d.tick();
    return 0;
}
\end{lstlisting}

\begin{takeawaybox}
Elevator design shines when you separate \emph{dispatch} (which car, a
\textbf{Strategy}) from \emph{motion} (how a car sweeps its stops, a
\textbf{State}/SCAN algorithm). Discuss the scheduling goal explicitly.
Follow-ups: concurrency across cars, priority/express calls, and
starvation avoidance.
\end{takeawaybox}

\section{Design a Vending Machine}
\label{sec:case-vending}

\problemstatement{Design a vending machine that holds products in slots,
accepts coins/notes, dispenses a selected product, returns change, and
handles out-of-stock and insufficient-funds cases.}

\begin{patternbox}[Clarify Before Designing]
Ask: What payment types (cash, card)? Should it make exact change, and
what if it cannot? Can a selection be cancelled mid-transaction? For this
design assume: cash only, refunds allowed, and a clear lifecycle from idle
to dispensing. The lifecycle is the heart of the problem.
\end{patternbox}

\subsection*{Identify the entities}

\texttt{VendingMachine} (the context/state machine), \texttt{Product},
\texttt{Slot}/\texttt{Inventory} (stock levels), \texttt{Coin}/payment
handling, and a set of \texttt{State} classes -- \texttt{IdleState},
\texttt{HasMoneyState}, \texttt{DispensingState}. The machine's behaviour
for the same action (``insert coin,'' ``select'') differs entirely by
state, which is the signal for the State pattern.

\begin{ideabox}[State Pattern Is the Backbone]
Model the machine as a \textbf{State} machine: each state handles
\texttt{insertCoin}, \texttt{selectProduct}, and \texttt{dispense}
appropriately and transitions the machine forward. This replaces a giant
\texttt{switch(currentState)} in every method and makes illegal
transitions (dispensing before paying) structurally impossible.
\end{ideabox}

\subsection*{Class design (C++ skeleton)}

\begin{lstlisting}
#include <bits/stdc++.h>
using namespace std;

class VendingMachine;
struct State {
    virtual void insertCoin(VendingMachine&, int amount) = 0;
    virtual void selectProduct(VendingMachine&, const string& code) = 0;
    virtual ~State() = default;
};

class VendingMachine {
    unique_ptr<State> state;
public:
    int balance = 0;
    unordered_map<string,int> stock;            // code -> quantity
    unordered_map<string,int> price;            // code -> price

    VendingMachine();                           // starts in IdleState
    void setState(unique_ptr<State> s) { state = move(s); }
    void insertCoin(int a) { state->insertCoin(*this, a); }
    void selectProduct(const string& c) { state->selectProduct(*this, c); }
};

struct IdleState : State {
    void insertCoin(VendingMachine& m, int a) override;
    void selectProduct(VendingMachine&, const string&) override {
        cout << "Insert money first\n";
    }
};
struct HasMoneyState : State {
    void insertCoin(VendingMachine& m, int a) override { m.balance += a; }
    void selectProduct(VendingMachine& m, const string& code) override {
        if (m.stock[code] == 0) { cout << "Out of stock\n"; return; }
        if (m.balance < m.price[code]) { cout << "Insufficient funds\n"; return; }
        m.stock[code]--;
        int change = m.balance - m.price[code];
        m.balance = 0;
        cout << "Dispensed " << code << ", change " << change << "\n";
        m.setState(make_unique<IdleState>());   // back to idle
    }
};

void IdleState::insertCoin(VendingMachine& m, int a) {
    m.balance += a;
    m.setState(make_unique<HasMoneyState>());   // transition
}
VendingMachine::VendingMachine() : state(make_unique<IdleState>()) {}

int main() {
    VendingMachine m;
    m.stock["A1"] = 2; m.price["A1"] = 30;
    m.selectProduct("A1");                      // "Insert money first"
    m.insertCoin(20); m.insertCoin(20);
    m.selectProduct("A1");                      // dispenses, change 10
    return 0;
}
\end{lstlisting}

\begin{takeawaybox}
The vending machine is the textbook \textbf{State} pattern problem:
recognise that the same operations behave differently per mode and give
each mode a class that acts and transitions. Follow-ups: card payments
(Strategy), exact-change detection, and cancel/refund as another
transition.
\end{takeawaybox}

\section{Design Tic-Tac-Toe}
\label{sec:case-tic-tac-toe}

\problemstatement{Design a Tic-Tac-Toe game (generalisable to an $n\times
n$ board) supporting two players taking turns, move validation, and
win/draw detection. Keep the design extensible to other grid games.}

\begin{patternbox}[Clarify Before Designing]
Ask: Fixed $3\times3$ or general $n\times n$ with $k$-in-a-row to win? Two
human players or human-vs-AI? For this design assume: general board size,
$n$-in-a-row to win, two players, and an efficient $O(1)$ win check per
move rather than rescanning the board.
\end{patternbox}

\subsection*{Identify the entities}

\texttt{Game} (orchestrates turns and end conditions), \texttt{Board}
(grid state and win detection), \texttt{Player} (has a mark/symbol),
\texttt{Move}, and \texttt{Piece}/\texttt{Symbol}. The game is a small
lifecycle: in-progress, won, drawn -- a light State/enum.

\begin{ideabox}[Incremental Win Check in O(1) per Move]
Instead of scanning all lines after every move, keep running counts per
row, per column, and the two diagonals for each player. A move updates
just its row, column, and (if on a diagonal) the diagonal counters; a
count reaching $n$ means a win. This turns win detection from $O(n^2)$
into $O(1)$ per move.
\end{ideabox}

\subsection*{Class design (C++ skeleton)}

\begin{lstlisting}
#include <bits/stdc++.h>
using namespace std;

class Board {
    int n;
    vector<vector<char>> grid;
    vector<int> rowCnt, colCnt;                 // signed counts per player
    int diag = 0, anti = 0;
public:
    Board(int size) : n(size), grid(size, vector<char>(size, '.')),
                      rowCnt(size, 0), colCnt(size, 0) {}

    bool inBounds(int r, int c) const { return r >= 0 && r < n && c >= 0 && c < n; }
    bool isFree(int r, int c) const { return grid[r][c] == '.'; }

    // returns true if this move wins. player = +1 (X) or -1 (O)
    bool place(int r, int c, char mark, int player) {
        grid[r][c] = mark;
        rowCnt[r] += player; colCnt[c] += player;
        if (r == c) diag += player;
        if (r + c == n - 1) anti += player;
        int win = n * player;
        return rowCnt[r] == win || colCnt[c] == win || diag == win || anti == win;
    }
};

struct Player { string name; char mark; int id; };  // id = +1 or -1

class Game {
    Board board; vector<Player> players; int turn = 0; int movesLeft;
public:
    Game(int n, Player a, Player b)
        : board(n), players{a, b}, movesLeft(n * n) {}

    string move(int r, int c) {
        Player& p = players[turn];
        if (!board.inBounds(r, c) || !board.isFree(r, c)) return "Invalid move";
        bool won = board.place(r, c, p.mark, p.id);
        movesLeft--;
        if (won) return p.name + " wins";
        if (movesLeft == 0) return "Draw";
        turn ^= 1;                              // next player
        return "ok";
    }
};

int main() {
    Game g(3, {"Alice",'X',+1}, {"Bob",'O',-1});
    cout << g.move(0,0) << "\n";                // X
    cout << g.move(1,0) << "\n";                // O
    cout << g.move(0,1) << "\n";                // X
    cout << g.move(1,1) << "\n";                // O
    cout << g.move(0,2) << "\n";                // X wins (top row)
    return 0;
}
\end{lstlisting}

\begin{takeawaybox}
Tic-Tac-Toe rewards two things: a clean separation of \texttt{Game}
(rules/turns) from \texttt{Board} (state/win check), and the $O(1)$
incremental win-detection insight. Generalise to $n\times n$ up front.
Follow-ups: an AI player via \textbf{Strategy} (minimax), and other grid
games (Connect Four) reusing the board abstraction.
\end{takeawaybox}

\section{Design Splitwise (Expense Sharing)}
\label{sec:case-splitwise}

\problemstatement{Design an expense-sharing app: users form groups, add
expenses split among members (equally, by exact amount, or by
percentage), and query who owes whom. Support simplifying debts so the
fewest transactions settle everyone.}

\begin{patternbox}[Clarify Before Designing]
Ask: What split types -- equal, exact, percentage, shares? Do we track
pairwise balances or a single net balance per user? Is debt simplification
required? For this design assume: all three split types, a net-balance
model, and an optional simplify step.
\end{patternbox}

\subsection*{Identify the entities}

\texttt{User}, \texttt{Group}, \texttt{Expense} (amount, payer,
participants, and a \texttt{Split}), a \texttt{SplitStrategy}
(\texttt{Equal}, \texttt{Exact}, \texttt{Percent}), and a
\texttt{BalanceSheet}/\texttt{ExpenseManager} that maintains net balances.
The varying part -- how an amount is divided -- is a textbook Strategy.

\begin{ideabox}[Split Logic Is a Strategy; Balances Are Netted]
Model each split rule as a \textbf{SplitStrategy} producing a
per-participant owed amount, so new rules never touch the manager. Store a
single net balance per (debtor, creditor) or per user; adding an expense
just adjusts these. Debt simplification is then a graph settle: net each
person, then greedily match biggest creditor with biggest debtor.
\end{ideabox}

\subsection*{Class design (C++ skeleton)}

\begin{lstlisting}
#include <bits/stdc++.h>
using namespace std;

struct SplitStrategy {                          // Strategy: division rule
    // returns each participant's share of `amount`
    virtual vector<double> split(double amount, int n,
                                 const vector<double>& params) const = 0;
    virtual ~SplitStrategy() = default;
};
struct EqualSplit : SplitStrategy {
    vector<double> split(double amount, int n, const vector<double>&) const override {
        return vector<double>(n, amount / n);
    }
};
struct PercentSplit : SplitStrategy {
    vector<double> split(double amount, int n, const vector<double>& pct) const override {
        vector<double> out(n);
        for (int i = 0; i < n; i++) out[i] = amount * pct[i] / 100.0;
        return out;
    }
};

class ExpenseManager {
    // balance[a][b] = how much a owes b
    map<string, map<string, double>> balance;
public:
    void addExpense(const string& payer, const vector<string>& members,
                    double amount, const SplitStrategy& strat,
                    const vector<double>& params = {}) {
        auto shares = strat.split(amount, members.size(), params);
        for (size_t i = 0; i < members.size(); i++) {
            if (members[i] == payer) continue;
            balance[members[i]][payer] += shares[i];   // member owes payer
        }
    }
    void showBalances() {
        for (auto& [debtor, m] : balance)
            for (auto& [creditor, amt] : m)
                if (amt > 1e-9)
                    cout << debtor << " owes " << creditor << ": " << amt << "\n";
    }
};

int main() {
    ExpenseManager mgr;
    vector<string> group = {"Alice", "Bob", "Carol"};
    mgr.addExpense("Alice", group, 300, EqualSplit{});   // each owes Alice 100
    mgr.showBalances();
    return 0;
}
\end{lstlisting}

\begin{takeawaybox}
Splitwise is a Strategy showcase: isolate the split rule so equal/exact/
percent are plug-ins, and keep a simple netted balance model. The
memorable extension is \emph{debt simplification} -- reduce the balance
graph to the minimum number of settling transactions by greedily matching
the largest creditor and debtor.
\end{takeawaybox}

\section{Design a Notification System}
\label{sec:case-notification}

\problemstatement{Design a system that sends notifications to users across
multiple channels (email, SMS, push), lets users set channel preferences,
supports templated messages, and can add new channels without disturbing
existing code.}

\begin{patternbox}[Clarify Before Designing]
Ask: Which channels, and can more be added later? Are messages templated?
Do users choose channels per notification type? Synchronous or queued? For
this design assume: pluggable channels, per-user preferences, templated
content, and asynchronous-friendly dispatch.
\end{patternbox}

\subsection*{Identify the entities}

A \texttt{NotificationChannel} interface (\texttt{Email}, \texttt{SMS},
\texttt{Push}), a \texttt{ChannelFactory} to create them, a
\texttt{NotificationService} that orchestrates sending, \texttt{User}
preferences, and optionally \texttt{Observer} so subscribers of an event
are notified. Message building can use a \textbf{Builder} or template
engine.

\begin{ideabox}[Channels Are Strategies Behind a Factory; Events Fan Out via Observer]
Each channel implements a common \texttt{send} interface (a
\textbf{Strategy}/plug-in); a \textbf{Factory} maps a channel name to its
implementation, so a new channel is one new class plus a registry entry.
For ``notify everyone interested in event X,'' layer \textbf{Observer} on
top: subscribers register, the event publishes, each is sent via their
preferred channel.
\end{ideabox}

\subsection*{Class design (C++ skeleton)}

\begin{lstlisting}
#include <bits/stdc++.h>
using namespace std;

struct NotificationChannel {                    // Strategy per channel
    virtual void send(const string& to, const string& msg) = 0;
    virtual ~NotificationChannel() = default;
};
struct EmailChannel : NotificationChannel {
    void send(const string& to, const string& msg) override {
        cout << "Email to " << to << ": " << msg << "\n";
    }
};
struct SmsChannel : NotificationChannel {
    void send(const string& to, const string& msg) override {
        cout << "SMS to " << to << ": " << msg << "\n";
    }
};

struct ChannelFactory {                         // Factory: name -> channel
    static unique_ptr<NotificationChannel> create(const string& type) {
        if (type == "email") return make_unique<EmailChannel>();
        if (type == "sms")   return make_unique<SmsChannel>();
        return nullptr;
    }
};

struct User { string id; vector<string> preferredChannels; string contact; };

class NotificationService {
public:
    void notify(const User& user, const string& message) {
        for (const auto& ch : user.preferredChannels) {
            auto channel = ChannelFactory::create(ch);
            if (channel) channel->send(user.contact, message);   // fan out
        }
    }
};

int main() {
    NotificationService service;
    User u{"u1", {"email", "sms"}, "alice@x.com"};
    service.notify(u, "Your order shipped");    // sent on both channels
    return 0;
}
\end{lstlisting}

\begin{takeawaybox}
A notification system is a clean \textbf{Strategy + Factory} design:
channels behind one interface, created by name, so new channels never
touch dispatch. Add \textbf{Observer} for event-driven fan-out and
\textbf{Decorator} for cross-cutting concerns (retry, rate-limit,
encryption) around a channel. Discuss queuing for reliability at scale.
\end{takeawaybox}

\section{Design a Logging Framework}
\label{sec:case-logger}

\problemstatement{Design a logging framework supporting severity levels
(DEBUG, INFO, WARN, ERROR), multiple output sinks (console, file,
network), a configurable minimum level, and the ability to add sinks or
formats without changing call sites.}

\begin{patternbox}[Clarify Before Designing]
Ask: What levels and can they be filtered globally or per sink? Multiple
simultaneous sinks? Thread-safe? Should a message be passed along a chain
of level handlers or broadcast to all sinks? For this design assume:
standard levels, multiple sinks, a global threshold, and thread safety as
a discussion point.
\end{patternbox}

\subsection*{Identify the entities}

A \texttt{Logger} (front door, often a \textbf{Singleton}), a
\texttt{LogLevel} enum, \texttt{LogSink}/\texttt{Appender} implementations
(\texttt{ConsoleSink}, \texttt{FileSink}), and a \texttt{Formatter}. Two
classic pattern framings apply: broadcast to all sinks (\textbf{Observer})
or pass up a \textbf{Chain of Responsibility} of level handlers.

\begin{ideabox}[Level Filtering + Pluggable Sinks]
Keep a single global threshold so sub-threshold messages are dropped
cheaply. Model each output as a \texttt{LogSink} plug-in the logger writes
to -- adding a sink never touches call sites. A \textbf{Chain of
Responsibility} variant instead routes a message through DEBUG$\to$INFO
$\to$WARN$\to$ERROR handlers, each deciding whether to act; both are valid,
name the trade-off.
\end{ideabox}

\subsection*{Class design (C++ skeleton)}

\begin{lstlisting}
#include <bits/stdc++.h>
using namespace std;

enum class Level { Debug, Info, Warn, Error };

struct LogSink {                                // pluggable output
    virtual void write(Level lvl, const string& msg) = 0;
    virtual ~LogSink() = default;
};
struct ConsoleSink : LogSink {
    void write(Level, const string& msg) override { cout << "[console] " << msg << "\n"; }
};
struct FileSink : LogSink {
    void write(Level, const string& msg) override { cout << "[file] " << msg << "\n"; }
};

class Logger {                                  // Singleton front door
    Level threshold = Level::Info;
    vector<unique_ptr<LogSink>> sinks;
    Logger() = default;
public:
    static Logger& instance() { static Logger l; return l; }
    void setThreshold(Level l) { threshold = l; }
    void addSink(unique_ptr<LogSink> s) { sinks.push_back(move(s)); }

    void log(Level lvl, const string& msg) {
        if (lvl < threshold) return;            // cheap filter
        for (auto& s : sinks) s->write(lvl, msg);   // broadcast to sinks
    }
};

int main() {
    Logger& log = Logger::instance();
    log.addSink(make_unique<ConsoleSink>());
    log.addSink(make_unique<FileSink>());
    log.setThreshold(Level::Info);
    log.log(Level::Debug, "verbose");           // filtered out
    log.log(Level::Error, "something broke");   // goes to both sinks
    return 0;
}
\end{lstlisting}

\begin{takeawaybox}
A logger design lets you name several patterns naturally: \textbf{Singleton}
for the global logger, pluggable \textbf{Sinks} (Strategy/Observer) for
outputs, and optionally a \textbf{Chain of Responsibility} over levels.
The must-mention non-functional concern is \emph{thread safety} -- a lock
or a lock-free queue around the sink writes at scale.
\end{takeawaybox}

\section{Design a Rate Limiter}
\label{sec:case-rate-limiter}

\problemstatement{Design a rate limiter that caps how many requests a
client may make in a time window -- e.g.\ 100 requests per minute -- and
lets the algorithm (token bucket, sliding window, fixed window) be swapped.
Return allow/deny per request.}

\begin{patternbox}[Clarify Before Designing]
Ask: Per-user, per-IP, or global limits? Which algorithm -- fixed window,
sliding window log/counter, token bucket, leaky bucket? Distributed across
servers or single-node? For this design assume: per-client keys, a
swappable algorithm, and single-node with a note on distributed
extensions.
\end{patternbox}

\subsection*{Identify the entities}

A \texttt{RateLimiter} facade keyed by client, and a
\texttt{RateLimitStrategy} interface with implementations
(\texttt{TokenBucket}, \texttt{FixedWindow}, \texttt{SlidingWindow}). Each
strategy stores per-client state and answers \texttt{allow()}. The varying
algorithm is, again, a Strategy.

\begin{ideabox}[Token Bucket in a Nutshell]
A token bucket holds up to \texttt{capacity} tokens, refilled at a steady
\texttt{rate} per second. Each request tries to take one token: succeed and
proceed, or fail and get rejected. It naturally allows short bursts (up to
capacity) while bounding the long-run average -- which is why it is the
most common choice. Refill lazily by computing elapsed time on each call
rather than running a timer.
\end{ideabox}

\subsection*{Class design (C++ skeleton)}

\begin{lstlisting}
#include <bits/stdc++.h>
using namespace std;

struct RateLimitStrategy {
    virtual bool allow(const string& client) = 0;
    virtual ~RateLimitStrategy() = default;
};

class TokenBucket : public RateLimitStrategy {
    double capacity, refillPerSec;
    struct Bucket { double tokens; double lastRefill; };
    unordered_map<string, Bucket> buckets;

    double now() {
        return chrono::duration<double>(
            chrono::steady_clock::now().time_since_epoch()).count();
    }
public:
    TokenBucket(double cap, double rate) : capacity(cap), refillPerSec(rate) {}
    bool allow(const string& client) override {
        double t = now();
        auto& b = buckets.try_emplace(client, Bucket{capacity, t}).first->second;
        b.tokens = min(capacity, b.tokens + (t - b.lastRefill) * refillPerSec);
        b.lastRefill = t;
        if (b.tokens >= 1.0) { b.tokens -= 1.0; return true; }   // consume
        return false;                                            // throttled
    }
};

class RateLimiter {                             // facade over the strategy
    unique_ptr<RateLimitStrategy> strategy;
public:
    RateLimiter(unique_ptr<RateLimitStrategy> s) : strategy(move(s)) {}
    bool allow(const string& client) { return strategy->allow(client); }
};

int main() {
    RateLimiter limiter(make_unique<TokenBucket>(3, 1));  // burst 3, 1/s
    for (int i = 0; i < 5; i++)
        cout << "req " << i << ": " << (limiter.allow("user1") ? "OK" : "DENY") << "\n";
    return 0;
}
\end{lstlisting}

\begin{takeawaybox}
Rate limiting is a \textbf{Strategy} problem: token bucket, fixed window,
and sliding window are interchangeable behind one \texttt{allow()}
interface. Know token bucket's burst-vs-average trade-off cold, and
compute refill lazily. The senior-level extension is \emph{distributed}
limiting -- shared counters in Redis with atomic increments and expiry.
\end{takeawaybox}

